On frappe en même temps un diapason A de \qty{\v{f1}}{Hz} et un diapason B dont la fréquence est inconnue. On entend \num{\v{fb}} battements par seconde.

\begin{enumerate}[label=\alph*)]
    \item Quelles sont les deux valeurs possibles de la fréquence du diapason B ?
    \item On colle un petit morceau de cire sur une branche du diapason B, ce qui diminue légèrement sa fréquence. La fréquence des battements augmente. Quelle était la fréquence du diapason B avant l'ajout de la cire ? Justifiez.
    \item Combien de battements entend-on en \qty{\v{t}}{s} (sans la cire) ?
\end{enumerate}
